% ICLR version, converted from the ACL ARR manuscript (Chekhov_paper_ACL/overrealization.tex).
% NOTE: the ICLR 2027 style files (iclr2027.zip) were not yet published at conversion
% time (2026-07-29); this uses the ICLR 2026 style files, which historically differ
% from year to year only in the year string. When iclr2027_conference.{sty,bst} are
% released at https://github.com/ICLR/Master-Template, drop them in and update the
% two \usepackage / \bibliographystyle lines below.
\documentclass{article}
\usepackage{iclr2026_conference,times}

\usepackage[T1]{fontenc}
\usepackage[utf8]{inputenc}
\usepackage{microtype}
\usepackage{inconsolata}
\usepackage{graphicx}
\usepackage{booktabs}
\usepackage{amsmath}
\usepackage{hyperref}
\usepackage{url}
\usepackage{todonotes}   % writing plan; remove \todo notes before submission

\title{When Chekhov's Gun Fires: Generation and Detection of\\Over-Realization in Text-to-Image Models}

% ICLR submissions are double blind: the author block below is hidden (replaced by
% "Anonymous authors") as long as \iclrfinalcopy remains commented out.
\author{The Chekhov Autonomous Research Pipeline \\
  \texttt{research/runs/20260724\_chekhovs-gun} \\}

%\iclrfinalcopy % Uncomment for camera-ready version, but NOT for submission.

\begin{document}
\maketitle

\begin{abstract}
% TODO(final): write last, after results are in.
% Skeleton of the arc:
% (1) Construct: content mentioned in the input is often *realized* in the output
%     even when the context does not warrant it -- over-realization.
% (2) Taxonomy: six families in three tiers of licensing -- space builders (1-4),
%     the use--mention distinction (5), and pragmatic relevance (6).
% (3) Dataset: paired minimal-difference items with positive controls, text-to-image.
% (4) Two tasks, not one: GENERATION (does the model suppress?) and DETECTION (can a
%     model tell, given the prompt, whether an image over-realized?).
% (5) Findings: which families fail; and that detection is not the easy half --
%     the same judge moves from kappa 0.23 to 0.94 on identical images on wording alone.
% (6) The text control: on identical stimuli four of six families never fail in text
%     and fail readily in images, which is what makes T2I the right setting.
\emph{Placeholder abstract.}
\emph{This manuscript was generated autonomously by an AI research pipeline; see the AI Use Statement.}
\end{abstract}

\section{Introduction}
\label{sec:intro}

\todo[inline]{\textbf{P1 --- Hook: the \emph{dramatic} framing, and why real inputs are
not dramatic.} Open with Chekhov's principle: a gun hung on the wall in act one must
fire in act three. Note that this is a principle about \emph{deliberate dramatic
setup}---the gun is a plant, placed by an author who intends it to pay off. Then the
turn: real user inputs are not dramatic. A passing remark, an ordinary topic shift, a
figure of speech, or a passage the retriever happened to return is not a plant and is
not foreshadowing. Yet generative models treat these incidental mentions as if they
were dramatic setups, and duly discharge them. Land this on Figure~\ref{fig:teaser}
(user will supply): (a) figurative---``a singer whose voice was like silk'' yields silk
fabric in the scene; (b) perspectival---``Anna reads a letter that says `I'm sorry'\,''
yields a letter lying flat to the camera, legible to the viewer but not to Anna.
Each with its paired control. Close the paragraph on the normative claim: faithful
generation frequently requires leaving the gun hanging.}

\todo[inline]{\textbf{P2 --- The construct: over-realization.} A mention is not a
licence to realize. Define \emph{over-realization}: content that appears in the input
but should not be realized in the output. The unifying observation is that
suppression is licensed by a range of devices, of which negation is only one---attitude
predicates, figurative frames, viewpoint frames, the use--mention distinction and plain
topical relevance all do the same work.
Introduce \emph{space builders} here, and be explicit that this is Fauconnier's
established term from Mental Spaces theory, not one we coin (the only term we coin is
\emph{over-realization}): a space builder opens a space distinct from the discourse
base space, and entities living in a non-base space should not be realized at base
level. State that this lets us treat a set of superficially unrelated defect classes as one
phenomenon with one measurement design.}

\todo[inline]{\textbf{P3 --- Position against prior work} (one short paragraph; detail
belongs in Section~\ref{sec:related}). Two gaps. (i) Existing work is concentrated on
\emph{negation} (NEG-TTOI, NEGATE, SpaceVLM, TNG-CLIP in vision; IFEval forbidden words
and NegE/Modality errors in text), with the remaining operators treated---if at
all---as separate defects with separate fixes, so nothing is comparable across them.
(ii) The recognition side is treated as solved: existing benchmarks assume a judge and
evaluate generators, whereas we find that judging over-realization is itself unreliable
enough to be worth measuring, and make it the second task.
Then \emph{one sentence} on contextual entrainment \citep{liu2026entrainment}: it
establishes that context-mentioned content is internally primed---probability-elevated
in the model's representations---and we take that as given; our question is the
complementary one of when that priming is actually \emph{realized} in output. Useful
framing if a sentence is needed: entrainment measures ``the word is on the tip of the
tongue''; we measure when it is said aloud. Keep this to one or two sentences so the
opening does not read as building on someone else's foundation.}

\todo[inline]{\textbf{P4 --- Why text-to-image, and why two tasks.} Two moves to make
here, both settled by the pilot and both needing one short paragraph.

\emph{Setting.} The benchmark is text-to-image. Justify it with the control rather than
by assertion: run in text, four of the six families never fail at all---negation,
attribution, figurative and use--mention sit at exactly $0.00$ across three models---while
the same stimuli fail readily as image prompts. What is marked \emph{lexically} is
respected; what has to be inferred from a viewpoint or a relevance relation is not, and
only the image modality exercises the second kind at scale. The text runs stay in the
paper as a control that establishes this asymmetry (\S\ref{sec:experiments}), not as half
a benchmark.

\emph{Two tasks.} Generation asks whether a model suppresses. Detection asks whether a
model can tell, given the prompt, that an image over-realized. Motivate the second from
our own experience rather than from principle: the same 7B judge moved from
$\kappa=0.23$ to $\kappa=0.94$ on identical images purely from how the question was
worded, so recognising over-realization is an open problem, not a solved preliminary that
can be delegated to an off-the-shelf judge.}

\todo[inline]{\textbf{P5--P6 --- deferred until experiments are complete.} P5: what we
find, with concrete numbers---family-wise ranking, where detection succeeds and fails,
and the text control. P6: the contributions and the data release statement.}

\begin{figure}[t]
\centering
\fbox{\parbox[c][3.2cm][c]{0.8\textwidth}{\centering\emph{Teaser figure ---
to be supplied by the author.}\\[2pt]\footnotesize (a) figurative: ``voice like silk''
$\rightarrow$ silk appears;\\ (b) perspectival: the letter faces the camera, not Anna.\\
Each with its paired control.}}
\caption{Two over-realization failures. In both cases the mention is incidental rather
than a dramatic plant, yet the model discharges it.}
\label{fig:teaser}
\end{figure}

\section{Related Work}
\label{sec:related}
% Planned content (~0.75 page). Four paragraphs, each ending in an explicit contrast.
% (1) Compositional / faithfulness evaluation of generation:
%     T2I-CompBench(++), TIFA, VQAScore, Commonsense-T2I (closest methodology:
%     paired adversarial prompts + MLLM judge), PhyBench, VBench-2.0.
% (2) Single-family prior work -- we adopt rather than claim these:
%     negation (NEG-TTOI, NEGATE, SpaceVLM, TNG-CLIP; IFEval forbidden words in text;
%     NegE/Modality/Tense errors in summarization taxonomies);
%     figurative (IRFL, T2I-ReasonBench idiom interpretation, SemEval-2025 ADMIRE,
%     I Spy a Metaphor, ViPE).
% (3) Adjacent but different in direction -- MUST be crisp:
%     POPE/CHAIR = the reverse direction (a VLM describing objects absent from an
%       image); we study input-mentioned content wrongly realized in output.
%     Theory-of-mind benchmarks (ToMi, FANToM, BigToM, OpenToM, Hi-ToM) test whether a
%       model can *reason about* beliefs; we test whether a model *suppresses*
%       embedded content when generating. Different task, different failure.
%     Typographic attacks (Goh et al. 2021 multimodal neurons; "Reading Isn't
%       Believing"; NAACL 2025 multi-image) are the DUAL of our family 5 (use--mention) and make a
%       good contrast: in recognition, text written in an image overrides what is
%       depicted ("reads first, looks later"); in generation, the word is rendered
%       instead of its referent. Text-rendering work (TextDiffuser, glyph-enhanced
%       generation) treats legibility as the goal and does not ask why the word is
%       there at all.
% (4) Theoretical sources: Fauconnier mental spaces / space builders; Genette
%     focalization and diegesis; Lakoff \& Johnson conceptual metaphor; Kaup's
%     two-step simulation of negation; Grice's maxim of relation.
\emph{Placeholder.}

\section{Construct and Taxonomy}
\label{sec:taxonomy}

\paragraph{Over-realization.}
We use \emph{over-realization} to denote content that is mentioned in the input but
should not be realized in the output. The term is deliberately placed within the
classical generation pipeline of \emph{content selection} followed by \emph{linguistic
realization}: an over-realization is content that should have been filtered during
selection but nonetheless reaches the surface. It is distinct from
\emph{over-specification} in referring expression generation (redundant attributes of
an already licensed referent), from \emph{overgeneration} (a generate-and-rank
technique, or a grammar accepting ill-formed strings), and from
\emph{over-literalization}, which we treat as one subtype within the figurative family.
We reserve \emph{intrusion} for the measured event (the mentioned content appearing in
the output) and \emph{base-space leakage} for the theoretical description below.

\todo[inline]{\textbf{The running example: one entity, six failures} --- and make this,
not the theory, carry the section. Present Table~\ref{tab:taxonomy} early and walk
through it: a single entity instantiated across all six families, so that the reader
sees the unity of the construct by demonstration rather than by argument. If the same
entity in the same kind of scenario can fail in six different ways, ``these are one
phenomenon'' needs no defending. This is also literally an instance of the benchmark's
core-pool design (Section~\ref{sec:dataset}), so the theoretical claim and the dataset
construction are the same move. \textbf{All six cells must be real generations from our
own runs, never invented} --- an example a reviewer cannot reproduce would discredit the
paper --- and which entity survives all six families is an empirical question to settle
in the pilot before this table is finalized.}

\todo[inline]{\textbf{Theory, compressed and attached to the examples.} Do not teach the
frameworks in the abstract; attach one sentence of theory to the example that motivates
it, so the citation reads as backing rather than as a lesson. For families 1--4, on the
negation row: in ``a room with no elephant'' the negation opens a space in which the
elephant exists, while the base space---what the picture should show---contains none;
\citet{fauconnier1994mental} calls such operators \emph{space builders}, and families
1--4 differ only in which space builder is used (negation and modality, attitude
predicates, figurative frames, viewpoint frames). Then two short extensions, each one
sentence: family 5 is licensed not by a space but by the use--mention distinction, since
the input contains a word form while licensing only its denotation; family 6 carries no
marker at all and is licensed pragmatically, by relation. \textbf{Calibrate the claim
carefully}: do not say that mental-spaces theory is incomplete or that we repair it---it
was never intended to cover relevance, for which Grice is canonical. Say instead that
\emph{for this construct}, space building is the dominant licensing mechanism but does
not exhaust them, and that we identify two further mechanisms. Cite Quine and Grice for
the extensions; the contribution we claim is the unification, not the discovery of those
mechanisms. Close by predicting what unification entails: if these really are one
phenomenon they should share a failure profile and respond to the same modulators, which
Section~\ref{sec:experiments} tests---so the taxonomy is load-bearing rather than
decorative. \textbf{Do not promise a representational mechanism here}: the
availability-versus-expression analysis was cut with the text half, and a prediction the
paper does not test is worse than no prediction.}

\todo[inline]{\textbf{$E$ is a different kind of thing in each family} (short, and
mostly carried by Table~\ref{tab:taxonomy}, which shows it rather than states it: the
same elephant is an object in family 1, a believed object in 2, a metaphor vehicle in 3,
occluded content in 4, a word form in 5, a topic in 6). Two consequences are worth one
sentence each. It is why the judges must be family-specific---detectors for objects, OCR
for word forms, structured geometric questions for viewpoint---which pre-empts the
question of why we do not use one uniform metric. And it splits family 4 in two:
\textbf{4a occlusion}, where $E$ is an object behind a barrier and can therefore use the
core pool and carry the aligned cross-family comparison, and \textbf{4b diegetic
legibility}, where $E$ is inscribed information read by an in-scene observer. 4b is the
most striking sub-type but cannot use the object pool, so it is a showcase rather than
part of the aligned comparison; note also that in 4b the content hardly matters---whether
the letter says ``I'm sorry'' or ``the meeting is at noon'' does not change whether it
faces the camera---so 4b varies the carrier (letter, phone screen, book, sign, chart)
instead.}

\todo[inline]{\textbf{What counts as realized, and two failure modes} (short, and
measurement-facing, since it determines the judges). The unit of realization is
\emph{information made accessible to the audience}, not merely an object being present:
this is what makes the perspectival family measurable, since the letter's content
belongs to Anna's viewpoint space and rendering it legible to the camera moves it into
the audience's. Then the two modes: families 1 and 4--6 are \emph{binary}---the question
is whether the content appears at all---while families 2 and 3 are \emph{marking}
failures, where the content may legitimately appear but only as explicitly embedded (a
thought bubble, or a hedge such as ``she believed''), and the failure is the loss of that
marker.}

\todo[inline]{\textbf{Why a linguistic account licenses an image benchmark} (still
needed: the taxonomy is built from linguistic devices but the benchmark scores pictures,
and a reviewer will ask why that is legitimate). The reason is simple and should be
stated plainly: \emph{the input is language in both cases}. A text prompt and an image prompt are
both linguistic, and mental-spaces theory is a theory of meaning construction from
linguistic input, so the licensing analysis is identical on the input side; only the
notion of realization differs on the output side, which the definition above supplies.
Note the precedent so this does not look like a stretch: cognitive-linguistic frameworks
have been extended to visual media for decades, notably in work on pictorial and
multimodal metaphor \citep{forceville1996pictorial}, where conceptual blending is used
directly for image analysis. \textbf{Then concede two genuine asymmetries}, here and
again in the Limitations. (i) Family 5 is not symmetric: in an image, any rendered
prompt glyph is unmistakably the signifier being realized, whereas a text output
consists of word forms throughout, so the text side is restricted to narrower cases such
as unfilled template slots and instruction echo. (ii) Family 2's marking is
grammaticalized in text (hedges and attribution verbs, precise and near-obligatory) but
merely conventional in images (thought bubbles, dream haze, framed insets---optional and
genre-bound), so marking judgements are cleaner in text and this is why the human
agreement numbers for that family will differ by modality.}

\begin{table}[t]
\centering
\small
\setlength{\tabcolsep}{4pt}
\begin{tabular}{p{2.1cm} p{1.8cm} p{4.9cm} p{4.0cm}}
\toprule
Family & Cue & Prompt ($S$, verbatim from the pilot) & Observed failure \\
\midrule
\textbf{1. Existence-canceling} & \emph{no, without}
 & ``a hotel lobby at midnight with \textbf{no} elephant in it''
 & An elephant stands in the lobby doorway \\
\addlinespace
\textbf{2. Attribution} & \emph{believes, thinks}
 & ``Maya at the ridge, \textbf{believing} that there is an elephant at the ridge''
 & An actual elephant on the ridge, unmarked \\
\addlinespace
\textbf{3. Figurative} & \emph{like, as \ldots\ as}
 & ``the porter arriving at the inn, \textbf{as heavy as} an elephant''
 & A literal elephant walks beside the porter \\
\addlinespace
\textbf{4a. Perspectival: occlusion} & \emph{on the far side, cannot see}
 & ``the view from where Tomas stands: a four-metre brick wall, with an elephant
   \textbf{entirely on the far side} of it''
 & The wall is split in two and the elephant shown in the gap \\
\addlinespace
\textbf{4b. Perspectival: legibility} & \emph{viewpoint frame}
 & ``Anna sits at the table \textbf{reading} a letter that says `I'm sorry'\,''
 & The letter faces the camera, legible to the viewer, not to Anna \\
\addlinespace
\textbf{5. Use--mention} & \emph{semiotic level}
 & ``a warehouse with a crate \textbf{of elephants}''
 & Predicted: the crate stamped \textsc{elephant}. Not observed on the piloted
   generator$^{\dagger}$ \\
\addlinespace
\textbf{6. Relevance} & \emph{none: pragmatic}
 & ``My cousin \ldots\ looks after an elephant \ldots\ \textbf{Anyway} --- a school
   gymnasium after the game.''
 & An elephant stands on the basketball court \\
\bottomrule
\end{tabular}
\caption{One entity, six failures, from real generations. Every prompt is a verbatim
$S$-condition stimulus from the pilot and every failure was observed (FLUX.1-dev), except
family 5$^{\dagger}$, where the piloted generator commits to no word form at all and the
predicted typographic failure did not occur; the cell awaits a generator with real text
rendering. The families are ordered by how explicitly the input cues the suppression:
families 1--4 carry an overt \emph{space builder} \citep{fauconnier1994mental}; family 5
turns on the use--mention distinction; family 6 carries no cue and must be inferred
pragmatically. Family 4 splits: 4a uses the shared entity pool and carries the aligned
cross-family comparison; 4b realizes \emph{information} rather than an object, uses a
carrier pool, and is the showcase. Definitions of each family are given in the running
text; this table's job is only to show that one entity fails six different ways.}
\label{tab:taxonomy}
\end{table}

\section{OverReal}
\label{sec:dataset}

\todo[inline]{\textbf{P1 --- Item structure: the S/P/A triple.} Every item is built
from a \emph{scenario} $Sc$ (the carrier situation), an \emph{entity} $E$ (the content
that may be over-realized) and the family's \emph{licensing device} $D$; the three
conditions are the same $(Sc,E)$ pair assembled the same ways, which is the clearest way
to make the design legible: $S = Sc + D(E)$ (``a room with no elephant''),
$P = Sc + E$ (``a room with an elephant''), $A = Sc$ (``a room''). $S$ is issued at two
strengths of cue --- explicit ($S_{\mathrm{exp}}$, the suppression stated) and implicit
($S_{\mathrm{imp}}$, the suppression left to inference) --- making explicitness a
measured within-family factor rather than only the ordering between families. $S$ asks whether the
model suppresses, $P$ whether it can still realize the same entity when licensed, $A$
what the coincidental base rate is. Say why $P$ is indispensable: without it a
degenerate model that realizes nothing scores perfectly. Define the headline metric
here---\emph{licence sensitivity} $\Delta = P(\text{realize}\mid P) -
P(\text{realize}\mid S)$, in $[-1,1]$, where 1 is perfect discrimination and 0 means the
device is ignored---and note it is the same paired effect size used throughout, so the
generation and detection halves share one quantity.}

\todo[inline]{\textbf{P1b --- The second task: detection.} Present it as a peer of
generation, not an appendix. Given a suppression prompt $S$ and a candidate image, decide
whether the image over-realizes $E$. Motivate it with the reliability result rather than
in the abstract: a 7B open VLM scored the \emph{same} images at $\kappa=0.23$ under a
three-way question and $\kappa=0.94$ under one plain binary, so the ability to recognise
over-realization is a capability to be measured, not a preliminary to be assumed.

The construction is what makes this non-circular, and it should be stated explicitly
because a reviewer will ask. Labels must not come from a model judge---that tests
detectors against a detector---and must not come from the generating condition either,
since an $S$ image may be a correct suppression and a $P$ image may have omitted $E$
altogether. Instead, \emph{pair prompts with images across conditions}, using the shared
seed that already makes $S$, $P$ and $A$ near-identical in composition:
(i)~$S$ prompt $+$ $P$ image $\rightarrow$ \textbf{positive}, labelled by construction ---
the image was generated from an explicit request for $E$;
(ii)~$S$ prompt $+$ $A$ image $\rightarrow$ \textbf{negative}, labelled by construction,
licensed by the measured base rate of $A$;
(iii)~$S$ prompt $+$ $S$ image $\rightarrow$ \textbf{human annotation}, the naturally
occurring distribution.

The first two classes are large and carry certain labels without anyone judging whether
the generator erred; only the third needs annotation. Report the
construction-guaranteed core and the human-annotated natural set separately---they
measure different things and averaging them hides that.}

\todo[inline]{\textbf{P1c --- Two controls the detection half cannot be published
without.} Run every detector twice more: once with the prompt withheld, once with a
\emph{mismatched} prompt. Above-chance performance in either condition means the items
leak---the detector is reading generator style, prompt length or composition artefacts
rather than the licensing relation---and the affected cells are void. Also state that
candidates come from at least two generators, so that recognising one model's style is
not a winning strategy, and that no evaluated detector may be a model whose labels were
used anywhere in construction.}

\todo[inline]{\textbf{P2 --- Evaluation protocol: how ``realized'' is decided.} This has
to come before the construction loop, which cannot select difficult items without a
score. Because $E$ is a different kind of thing in each family, the judges are
family-specific and should be presented as such rather than apologised for: presence for
families 1, 3, 4a and 6; the rendered word form for family 5; geometry---can the in-scene
observer perceive the depicted information---for family 4b.

\textbf{Report the protocol finding as a finding, because it is one.} Our first protocol
used multiple-choice questions with a third ``hedged'' option and reached
$\kappa=0.80/0.42/0.23$ against hand annotation on families 2, 3 and 4b. Replacing each
with a single positive binary question about depicted content reached
$\kappa=0.94/1.00/0.94$ on the same images, and $0.955$ and $1.00$ on two families
constructed afterwards and therefore out of sample. Four wording faults account for the
gap and are worth a short list, since they generalise beyond this benchmark: a third
hedged option gives an uncertain judge somewhere to park; listing the embedding devices
(``a picture, a statue, a logo\ldots'') makes the judge answer about the \emph{medium},
because every generated image is a picture; the words ``real''/``real live'' make it
answer about ontological status---asked whether a ``real live elephant'' was present it
answered \emph{no} on all 72 images, several dominated by an elephant; and negatively
framed questions are simply agreed with, at an 89\% yes-rate against a 42\% truth rate.

Note also that \emph{paired forced choice}, which we had planned as the primary protocol
on reliability grounds, cannot serve as one here: it measures discriminability rather
than realization, and it degenerates exactly where the phenomenon is, since an
over-realized $S$ output is genuinely indistinguishable from its $P$ counterpart. In our
pilot 7--9 of 12 pairs per family carried no information for this reason. Use absolute
per-image judgement for the rate, and paired forced choice only as a validity check on
the judge.}

\todo[inline]{\textbf{P3 --- Composition: crossed entities and scenarios.} Entity and
scenario are \emph{crossed random factors}, not nested: each entity appears in several
scenarios and each scenario hosts several entities, assigned by Latin-square
counterbalancing so no $(E,Sc)$ pair repeats. Nesting would confound the two---a failure
could be a sticky entity or a permissive scenario---while full crossing is impossible
because many combinations are unnatural. The payoff is that the analysis can fit
\texttt{failure \textasciitilde{} family $\times$ openness + (1|entity) + (1|scenario)}
and separate entity from scenario variance, which most benchmarks cannot do. The design
presupposes that $E$ would not arise in $Sc$ unprompted, so the $A$ condition is run
first and any pair whose base rate exceeds a threshold is dropped---$A$ therefore serves
triple duty: validity check, pair filter, and the source of certain negatives for the
detection half. Entities come from a \emph{small core pool} instantiated across the
families that carry the aligned comparison, plus family-specific pools chosen for what
each family actually tests. \textbf{Do not force one pool on every family}: the pilot
found essentially no entity variance---four candidate entities failed exactly the same
families---so a shared pool buys protection against a confound that is not there, while
costing each family the entities that make it expressive. Family 5 in particular needs
carriers that must bear text, and 4b needs carriers rather than contents. Because the
load-bearing factor differs by family, the entity-to-scenario ratio is family-specific:
20 conversations $\times$ 3 entities for relevance, where whether a mention reads as
incidental depends on the surrounding talk and not on which animal was named, and 20
carriers $\times$ 3 contents for 4b. Sixty items per family cell gives 360 combinations
and, with four conditions at three openness levels, 4{,}320 generation prompts; give
Commonsense-T2I ($\approx$150 pairs) and T2I-CompBench (8{,}000 prompts) as scale
reference points. Scenarios are stratified over domain, the syntactic position of the
licensing device, and low prior compatibility between $E$ and $Sc$; they are seeded from
real corpora and completed by hand where a stratum is thin.}

\todo[inline]{\textbf{P3b --- Naturalness is a reported property, not a filter.} State
that items are rated for naturalness on a human-annotated sample and that the rating is
reported next to difficulty, so the trade-off is visible rather than buried in
construction. Two things to be explicit about. (i) An item that is hard because it is odd
is weak evidence: our own first figurative items were built by exhaustive crossing and
produced strings like \emph{``the auditor entering the office, as hungry as a wolf''},
which no user would write. (ii) We tried and rejected the obvious automatic instrument:
per-token perplexity under a base LM is dominated by the predictability of the scenario
phrase and ranked our most awkward figurative item among the most natural, so it is not a
naturalness measure and must not be used as a floor inside the difficulty loop. Say
plainly that this is why naturalness is human-rated here.}

\todo[inline]{\textbf{P4 --- The openness ladder.} Each item is issued at three
openness levels: L1 a fully specified scene prompt (subject, setting, composition,
style), L2 a conventional prompt, L3 a sparse creative prompt. Openness is a designed ordinal factor, not a measured quantity---L1 through L3
are constructed to be increasingly free, and we do not assign them calibrated
values---so the claims it supports are within-modality ones: whether failure rises with
openness inside each modality, and whether the most constrained image setting already
exceeds the most open text setting.}

\todo[inline]{\textbf{P5 --- Construction as a difficulty loop.} A first pass of
templated items will not be challenging, so construction runs model-in-the-loop, with
precedent in adversarial dataset building \citep{nie2020anli,kiela2021dynabench,%
lebras2020aflite}. Each round: generate candidates, run a set of \emph{prober} models,
then mutate the items every prober passes (strengthen the marker, lengthen the distance,
choose a more tempting entity), inspect the items every prober fails---usually
construction bugs, such as an entity that turns out to be plausible in the scenario---and
keep the discriminative middle; two or three rounds. Every item ships with
\texttt{difficulty}, the fraction of probers that fail it, which supports reporting by
stratum and keeps the benchmark useful as models improve. State the bias plainly:
adversarial filtering overfits items to the probers, so prober and evaluated model sets
are disjoint (probe with one family, evaluate with another, and cross-validate by
swapping), and the pre- and post-filter subsets are released separately so the effect of
filtering is visible.}

\todo[inline]{\textbf{P6 --- Validation.} Automatic checks: no target string may occur
in the $P$ or $A$ context or in any instruction, and $S$ and $P$ are matched for length
and framing. Human checks on a stratified subset: that suppression is genuinely required
(contestable for families 2, 3 and 4) and judge--human agreement reported \emph{per
family and per modality}, since presence checks are reliable while the marking judgement
in family 2 and the geometric judgement in 4b are not---and, as noted in
Section~\ref{sec:taxonomy}, marking is grammaticalized in text but only conventional in
images, so agreement should be expected to differ by modality. Conclusions for those
families are only as strong as the agreement reported here.}

\begin{table}[t]
\centering
\small
\begin{tabular}{lccc}
\toprule
Family & Items & Gen prompts & Det items \\
\midrule
1. Existence-canceling & --- & --- & --- \\
2. Attribution         & --- & --- & --- \\
3. Figurative          & --- & --- & --- \\
4a. Perspectival (occlusion) & --- & --- & --- \\
4b. Perspectival (legibility) & --- & --- & --- \\
5. Use--mention        & --- & --- & --- \\
6. Relevance           & --- & --- & --- \\
\midrule
Total & --- & --- & --- \\
\bottomrule
\end{tabular}
\caption{OverReal statistics (placeholder). Items are scenario--entity
combinations; each contributes four conditions ($S_{\mathrm{imp}}$, $S_{\mathrm{exp}}$,
$P$, $A$) at three openness levels. Det items pair $S$ prompts with cross-condition
images (\S\ref{sec:dataset}).}
\label{tab:datastats}
\end{table}

\section{Experiments and Analysis}
\label{sec:experiments}

\subsection{Setup}
% Planned: generators (>= 2 open T2I models; FLUX.1-dev piloted, plus one with strong
% text rendering for family 5), detectors for the detection half, family-specific
% binary judges validated against human annotation, and licence sensitivity from S/P.

\todo[inline]{\textbf{On openness.} The three openness levels are a designed factor
rather than a measured one, so claims stay within the modality: failure rises with
openness across L1--L3. The cross-modal calibration problem that used to live here is
gone with the cross-modal design; what remains in the Limitations is that the levels are
constructed rather than calibrated.}

\subsection{Main results: which families fail}
% Planned: Table 3 = failure rate by family x generator x openness.
% Pilot expectation: perspectival (4a, 4b) and relevance are the weakest; use-mention
% did not reproduce on FLUX at all and needs a strong text-rendering generator.
\emph{Placeholder.}

\subsection{Detection: can a model recognise over-realization?}
% Planned: accuracy on the construction-guaranteed core and on the human-annotated
% natural set, reported separately; the prompt-withheld and mismatched-prompt controls
% reported alongside every number, not in an appendix.
\todo[inline]{\textbf{What this subsection must not do.} Do not report a single pooled
accuracy. The construction-guaranteed items and the human-annotated $S$ items measure
different things: the first asks whether a detector can relate a prompt to an image at
all, the second whether it can do so on the distribution the generator actually
produces. Report both, and report the two leakage controls beside them---any cell where
the prompt-withheld or mismatched-prompt condition beats chance is void and should be
printed as void rather than quietly dropped.}

\subsection{The text control: why this is a text-to-image benchmark}
\todo[inline]{\textbf{One subsection, one table, no more.} Run the same items as text
prompts on open LLMs and show the asymmetry that motivates the setting: four of the six
families never fail in text---negation, attribution, figurative and use--mention at
$0.00$---while the same content fails readily as image prompts. Only the perspectival
families and relevance produce text failures at all. State the generalisation plainly:
suppression that is marked \emph{lexically} is respected, suppression that must be
inferred from a viewpoint or a relevance relation is not. Then two caveats that must not
be omitted. (i) The text figures are for one item difficulty; these are single-clause
items with the device in the same sentence as the entity, and the adversarial
construction loop was never run on the text side. (ii) Scope-aware scoring is essential
and easy to get wrong---a plain string match rates the figurative family at up to $0.42$
when its true literal-realization rate is $0.00$, because ``like a tiger'' contains
``tiger''. Any figurative result scored without scope will manufacture an effect that is
not there.}

\section{Conclusion}
\label{sec:conclusion}
\emph{Placeholder.}

% ICLR note: unlike ACL ARR, a Limitations section is NOT exempt from the page limit.
% It is kept in the main text for now; move (parts of) it to the appendix if the
% 9-page budget gets tight.
\section*{Limitations}
\label{sec:limitations}
\paragraph{Language.} All stimuli are in English. Space builders are realized very
differently across languages---tense and aspect are morphologically marked in some
languages and absent in others, evidentiality is grammaticalized in some but lexical in
English, and negation scope varies---so neither the failure rates nor their ordering
across families should be assumed to transfer. A multilingual instantiation is the most
important extension of this work.

\paragraph{Model coverage.} We evaluate open-weight models only, so we cannot speak to
proprietary systems, and our image models span a limited range of architectures.

\paragraph{Judge reliability, and its sensitivity to question form.} Every automatic
judgement here is reliable only in the form we validated it in. The same 7B VLM scored
identical images at $\kappa=0.23$ and $\kappa=0.94$ depending on the wording of the
question, so our numbers should be read as properties of a (model, question) pair rather
than of the model. We report human agreement on a stratified subset for every family, and
release the exact question text; a re-implementation that rewords the questions is not
running our protocol.

\paragraph{Oblique realization.} We label, but do not automatically score, cases where
the content enters as a depiction the scene independently licenses---a drawing of a tiger
on a waiting-room wall, a calendar, an anthropomorphic blend. These are intrusions by our
definition of realization (information made accessible to the audience) and successes by
a narrower referential reading. Reliably separating ``a picture of a tiger'' from ``a
tiger'' proved harder than the rest of the protocol combined, so the category is
human-annotated and reported, not folded into the headline rate.

\paragraph{Detection labels.} Most detection items are labelled by construction rather
than by annotation, which is what keeps the task non-circular, but it also means their
difficulty is not matched to the naturally occurring distribution. The human-annotated
subset is the one that speaks to real difficulty, and it is smaller.

\paragraph{Stimulus naturalness.} Items are constructed from templates in order to
support minimal-difference pairing; naturally occurring contexts may differ in
distribution.

\paragraph{Scope.} The benchmark is text-to-image; video is covered only as a case study
(Appendix). We characterise \emph{when} over-realization happens and \emph{whether} it can
be detected, and offer no representational account of \emph{why}: the availability-versus-
expression analysis lives in the text modality, which we retain only as a control.

\paragraph{Text is a control, at one difficulty.} The asymmetry that motivates the T2I
setting---four of six families never failing in text---is established on single-clause
items with the licensing device in the same sentence as the entity, and the adversarial
construction loop was never run on the text side. It supports ``these families are solved
in text at this difficulty'', not ``these families are solved in text''.

% ICLR 2027: optional, at the end of the main text before references; does not count
% toward the page limit; at most 1 page.
\section*{Ethics Statement}
% Planned: synthetic stimuli, no personal data; the failure studied is a faithfulness
% risk rather than a safety hazard; mitigation-oriented rather than attack-oriented.
\emph{Placeholder.}

% ICLR 2027: recommended, paragraph-long, at the end of the main text before
% references; does not count toward the page limit. Reference (do not restate) the
% parts of the paper/appendix that support reproduction: released items and exact
% judge question text, generator/judge model identities, construction-loop protocol.
\section*{Reproducibility Statement}
\emph{Placeholder.}

% ICLR 2027: REQUIRED (see the ICLR AI Policy for Authors); does not count toward the
% page limit. This replaces the ACL "Disclosure" section.
\section*{AI Use Statement}
\label{sec:disclosure}
This manuscript---including its hypotheses, experimental design, code, analysis,
figures, and text---was generated autonomously by an AI research pipeline operating
without human-written experimental code and without external model APIs. Human
involvement was limited to specifying the research direction and compute constraints.

\nocite{*}  % skeleton stage: list all available references
\bibliography{references}
\bibliographystyle{iclr2026_conference}

\appendix

\section{Video Case Study}
\label{app:video}
\emph{Placeholder.}

\section{Full Prompts and Annotation Guidelines}
\label{app:prompts}
\emph{Placeholder.}

\end{document}
